\documentclass[10pt,a4paper]{amsart}
\usepackage[margin=19mm]{geometry}
\usepackage{amsmath,amssymb,mathtools}
\usepackage[scr=rsfs,cal=euler]{mathalfa}
\usepackage{xfrac}
\usepackage[unicode,hidelinks,bookmarks=false]{hyperref}
\newcommand{\ie}{i.e.\ }
\newcommand{\cf}{cf.\ }
\newcommand{\resp}{respectively\ }
\newcommand{\Subsection}{\subsection}
\providecommand{\nobreakdash}{\nobreak\hbox{-}}
\setlength{\emergencystretch}{2em}
\multlinegap0pt
\usepackage{bm}
\usepackage{bbm}
\usepackage{dsfont}
\usepackage{xspace}
\usepackage{cancel}
\usepackage{mathtools}
\usepackage{cleveref}
\usepackage{amsthm}
\makeatletter
\@ifundefined{theorem}{\newtheorem{theorem}{Theorem}}{}
\@ifundefined{lemma}{\newtheorem{lemma}[theorem]{Lemma}}{}
\makeatother
\newcommand{\tr}{\mathrm{tr}}
\title[Positive Moments Forever: Undecidable and Decidable Cases]{Positive Moments Forever: Undecidable and Decidable Cases}
\author{Gemma De les Coves, Joshua Graf, Andreas Klingler, Tim Netzer}
\date{Source: arXiv:2404.15053v2 (CC BY 4.0)}
\begin{document}
\maketitle

\noindent\emph{Source and licence.} Positive Moments Forever: Undecidable and Decidable Cases. Version arXiv:2404.15053v2 (CC BY 4.0), distributed under the Creative Commons Attribution 4.0 International licence, as recorded in the arXiv licence metadata for this exact version and shown on the abstract page https://arxiv.org/abs/2404.15053v2. Full rights notice: licenses/arxiv-2404.15053-CC-BY-4.0.txt.\par\smallskip

\noindent\textit{Editorial scope.} Counted unit: the Theorem whose source label is \texttt{thm:moment\_orthogonal} in the article (source0.tex lines 424--427, source label \texttt{thm:moment\_orthogonal}), delivered together with the article's own complete original proof of it (source0.tex lines 428--447, one \texttt{proof} environment of 20 lines). No mathematics is written, repaired or re-derived: statement and proof are the article's own text line for line. Required context retained uncounted: lem:compactAlgebraic (lines 392--396). Every definition, display and label of those lines is retained unchanged. Omitted from this package and not redistributed: 1--391, 397--423, 448--758. Nothing inside a retained excerpt is omitted or altered: every retained body is delivered exactly as the article's lines are written, so no omission receipt of any kind accompanies this package. Citations: the retained excerpts carry no citation command at all, so no bibliography entry of the article is transcribed and no reference is redistributed. Reference adaptation: none was needed and none was made; every reference inside the delivered unit resolves inside it, so no unresolved reference or citation is produced. Printed slips kept literal: no printed word, symbol, hypothesis or number of the counted statement or of its proof was altered. Layout and class: the article's own class is \texttt{article} with packages including graphicx, geometry, amssymb, amsmath, mathtools, enumitem, amsthm, hyperref, cleveref, caption, subcaption, authblk, microtype, natbib; the wrapper is the lane's pinned \texttt{amsart} editorial wrapper at 10pt on A4, which restates the theorem-like environments the retained text needs and adds the article's own macro definitions that the retained text needs (\texttt{\textbackslash tr}), with extra wrapper packages (bm, bbm, dsfont, xspace, cancel, mathtools, cleveref). The delivered numbering is the unit's own. Assets: no retained span contains a figure environment, a graphics-inclusion command, a PDF-page inclusion, a table or a TikZ picture; no image file is copied into this package, the package declares no asset dependency and no image file is redistributed with it. Licence: version arXiv:2404.15053v2 of this article is distributed under the Creative Commons Attribution 4.0 International licence, as recorded in the arXiv licence metadata for this exact version and shown on the abstract page https://arxiv.org/abs/2404.15053v2. A full rights notice is delivered at licenses/arxiv-2404.15053-CC-BY-4.0.txt. Characters: every retained excerpt is pure ASCII, so the delivered engine, XeLaTeX with its default Latin Modern text font, reports no missing character.
\begin{lemma}\label{lem:compactAlgebraic}
Let $A_1, \ldots, A_d \in \mathrm{O}_s(\mathbb{Q})$ be orthogonal $s \times s$ matrices with rational entries. Then $\mathcal{G} \coloneqq \overline{\langle A_1, \ldots, A_d \rangle}$ is a compact algebraic group.
Moreover there is a recursively enumerable sequence of rational polynomials $(p_k)_{k \in \mathbb{N}}$ defining $\mathcal{G}$ inside $\mathrm{Mat}_s(\mathbb R)$.

\end{lemma}

\begin{theorem}
\label{thm:moment_orthogonal}
The moment positivity problem for  $\mathcal D=\mathrm{O}_s(\mathbb{Q})$  is decidable.
\end{theorem}

\begin{proof}
We will present two procedures, each certifying either yes- or no-instances in finite time. Letting these algorithms run in parallel will result in a decision algorithm for the problem.

Certifying no-instances for $A \in \mathrm{O}_s(\mathbb{Q})$ is achieved by iteratively checking whether $\tr(A^n) \geqslant 0$ holds for every  $n$. If $A$ is a no-instance, this algorithm will halt when detecting $\tr(A^n) < 0$ for the first time.

We now present an algorithm to certify yes-instances in finite time. For a given $A \in \mathrm{O}_s(\mathbb{Q})$, the moment membership problem can be rephrased  as
$$\forall B \in \langle A \rangle: \tr(B) \geqslant 0.$$
By the continuity of the trace, this is equivalent to 
\begin{equation}
\label{eq:firstOrderStatement}
\forall B \in \overline{\langle A \rangle}: \tr(B) \geqslant 0.
\end{equation}
By \cref{lem:compactAlgebraic} there exists a recursively enumerable sequence of polynomials $(p_k)_{k \in \mathbb{N}}$ and some $n \in \mathbb{N}$ such that
$$\overline{\langle A\rangle} = \mathcal{V}(p_1,\ldots, p_n).$$
Now step $k$ of the algorithm verifies the statement
\begin{equation}\label{eq:VarietyStatement}
\forall B \in \mathcal{V}(p_1, \ldots, p_k)\colon \tr(B) \geqslant 0
\end{equation}
which is decidable by the Tarski--Seidenberg Theorem, since it is a statement in first order logic. As soon  as  \cref{eq:VarietyStatement} is true for the first time, the algorithm halts and outputs a correct yes-answer. This  will be the case after at most $n$ steps, if $A$ is a yes-instance.
\end{proof}

\end{document}
